\chapter{구현}
\label{ch:impl}

제 \ref{ch:method}, \ref{ch:error}장의 이론을 PyTorch로 구현하였다. 본 장은 그 과정에서
필요했던 설계 결정을 기술한다. 이들은 원 제안서에 없던 것이며, 이론을 실제로 실행
가능한 시스템으로 바꾸는 데 필수적이었으므로 그 자체로 보고할 가치가 있다.

\section{구성}

\begin{table}[htbp]
\centering
\small
\caption{모듈 구성과 대응하는 이론.}
\label{tab:modules}
\begin{tabular}{@{}llp{6.6cm}@{}}
\toprule
모듈 & 대응 & 내용 \\
\midrule
\texttt{core/}        & --- & 격자 기하, 설정, 결정성, 단계별 계측 \\
\texttt{levelset/}    & 제 \ref{sec:spacing-aware}--\ref{sec:warmstart}절
                      & 간격 인지 연산자, Chan--Vese, 부호 거리 도구, marching tetrahedra \\
\texttt{surfel/}      & 제 \ref{sec:canonical}--\ref{sec:transport}절
                      & 정준 집합, 법선 정사영, 접선 전송, 밀도 제어 \\
\texttt{render/}      & 제 \ref{sec:raster}절
                      & 2DGS 래스터라이저, thin-3DGS 및 mesh 기준선, 광선 행진 참조 \\
\texttt{losses.py}    & 식 \eqref{eq:total-loss} & 외형/mask/normal/distortion 항 \\
\texttt{residual/}    & 제 \ref{sec:residual}절 & 정칙화 최소제곱(CG), 저계수 압축 \\
\texttt{data/}        & 제 \ref{sec:phantom}절 & 합성 팬텀, ACDC/M\&Ms-2 로더 \\
\texttt{metrics/}     & 제 \ref{sec:metrics}절 & 모든 보고 지표 \\
\texttt{pipeline/}    & 제 \ref{sec:pipeline}절 & 전처리, 정준 fitting, 재생, 직렬화 \\
\texttt{experiments/} & 제 \ref{ch:protocol}장 & 이론 검증, RQ1--RQ6, ablation \\
\bottomrule
\end{tabular}
\end{table}

전체 49개 모듈, 약 13{,}900행이다.

\section{설계 결정}
\label{sec:impl-decisions}

\subsection{기하는 최적화 불가능하다}

surfel의 anchor, 접선축, 법선을 \texttt{nn.Module}의 \emph{버퍼}로 등록하고 파라미터로
등록하지 않았다. 따라서 어떤 손실도 이들에 도달할 수 없다. 최적화되는 것은 진폭,
불투명도, 두 접선 scale뿐이다.

이는 편의가 아니라 주장의 정합성 문제이다. 제 \ref{sec:raster}절에서 언급하였듯
기하가 학습 가능하면 본 방법은 조용히 자유형 동적 2DGS fitting이 되고, ``기하가
$\surft$에서만 온다''는 주장이 무효가 된다. 학습률을 0으로 두는 것으로도 같은 효과를
낼 수 있으나, 설정 실수 한 번으로 주장이 깨진다. 버퍼로 등록하면 구조적으로 불가능해진다.
구현의 테스트는 \texttt{anchor.requires\_grad}가 참이면 실패한다.

\subsection{감독 신호가 무엇인지 정의해야 했다}

원 제안서는 감독을 ``알려진 slice plane의 mask, depth, normal, intensity''로 제한한다.
그러나 cine CMR은 복셀 볼륨이고 사진이 아니므로, 이 문장을 실행 가능한 목표로
바꾸어야 했다.

채택한 정의는 다음이다. 알려진 slice plane 카메라의 각 화소 광선에 대해 저장된
$\surft$와의 첫 교차를 광선 행진으로 찾고, 그 지점에서
\begin{itemize}[leftmargin=2em]
  \item 교차 여부 $\to$ 참조 silhouette, $\mathcal{L}_{\mathrm{mask}}$의 목표,
  \item 교차까지의 거리 $\to$ 참조 $\depthq$, depth RMSE의 기준,
  \item $\gradh\levt/\norm{\gradh\levt}$ $\to$ 참조 법선, 명제
  \ref{prop:normal-error}의 각오차 기준,
  \item 그 지점의 MRI 값 $\to$ $\mathcal{L}_{\mathrm{app}}$의 목표
\end{itemize}
를 읽는다. 이 정의의 두 성질이 중요하다. 첫째, 네 신호 모두 level-set과 MRI에서만
오므로 fitting이 저장 표현에 없는 정보를 받지 않는다. 둘째, mesh, thin-3DGS, 2DGS가
\emph{동일한} 참조에 대해 채점되므로 RQ5/RQ6가 표현을 비교하며 분할을 비교하지 않는다.

\subsection{좁은 띠는 희소 목록이 아니라 조밀 crop이다}

교과서적 좁은 띠는 활성 복셀의 희소 목록을 유지한다. GPU에서는 이것이 작은 상자에
대한 조밀 연산보다 느린 경우가 많다. 따라서 띠를
$\{\abs{\lev}<b\}$의 경계 상자에 유한차분 halo를 더한 \emph{조밀 crop}과, 그 안에서의
masked 갱신으로 실현하였다. 속도 이득은 crop에서 나온다.

한 가지 함정이 있다. 식 \eqref{eq:cin}, \eqref{eq:cout}의 영역 평균은 $\dom$ 전체에
대한 적분이다. crop 위에서 계산하면 외부 영역이 얇은 껍질로 축소되어 $\cout$이
편향된다. 따라서 영역 평균만은 \emph{전체} 볼륨에서 다시 계산한다. 이 비용은 두 번의
축약이므로 곡률 계산에 비해 무시할 만하다.

\subsection{빠른 래스터라이저의 두 근사}

성능을 위한 타일 기반 래스터라이저는 식 \eqref{eq:compositing}에 두 가지 근사를
도입한다.

\begin{enumerate}[leftmargin=2em]
  \item 타일 내 정렬은 화소별 정확한 $\depthq^t_i(q)$가 아니라 anchor 깊이로 한다.
  \item 타일당 가장 가까운 $M$개만 유지하고 나머지를 버린다.
\end{enumerate}

이 두 근사가 무해하다고 \emph{가정하지 않기 위해}, 모든 surfel을 모든 화소에 대해
평가하고 화소별로 정확한 $\depthq$로 정렬하는 브루트포스 참조 래스터라이저를 함께
구현하였다. 작은 영상에서 두 구현의 차이를 측정할 수 있으며, 구현의 테스트는 잘
분리된 surfel 배치에서 두 결과가 $10^{-3}$ 이내로 일치할 것을 요구한다. 절단된 페어
수와 절단율도 매 렌더링에서 보고한다.

\subsection{Depth distortion의 대체}

원 2DGS의 depth distortion 항은 $\sum_i\sum_j w_i w_j\abs{\depthq_i-\depthq_j}$로
$O(M^2)$이다. 구현은 $O(M)$인 가중 분산
\begin{equation}
  \sum_i w_i\bigl(\depthq_i - \bar\depthq\bigr)^2,
  \qquad \bar\depthq = \frac{\sum_i w_i \depthq_i}{\sum_i w_i}
  \label{eq:distortion-surrogate}
\end{equation}
로 대체하였다. 두 형태 모두 한 광선의 모든 기여가 한 깊이에 모일 때 정확히 소멸하므로
정규화의 목적은 보존된다. 이는 참조 구현으로부터의 의도적 이탈이므로 명시한다.

\subsection{임의 시점 탐색 문제}
\label{sec:seek-problem}

식 \eqref{eq:store}은 프레임별 anchor를 저장하지 않는다. 따라서 저장된 것만으로
프레임 $t$를 표시하려면 \emph{정준} anchor를 $\surft$로 직접 정사영해야 한다. 그러나
전처리는 anchor를 $\anchor^{t-1}\to\anchor^t$로 \emph{순차} 정사영하였다. 두 경로는
동일하지 않다. 순차 정사영은 경로를 누적하지만 직접 정사영은 훨씬 긴 한 걸음을 밟으므로
trust region에 걸리거나 잘못된 표면 분기에 착지할 가능성이 크다.

어느 한쪽을 조용히 택하는 대신 구현은 두 방식을 모두 제공하고 그 격차를 측정하는
함수를 제공한다. 직접 정사영이 $\Esurf$를 유의하게 악화시킨다면, 정직한 저장량 회계는
프레임별 anchor를 추가하거나 viewer를 순차 재생으로 제한해야 한다. 이는 표현에 대한
\emph{발견}이며 버그가 아니다.

\caveat{%
이 문제는 원 제안서에서 다루어지지 않았다. 구현 과정에서 드러난 것이며, 식
\eqref{eq:store}의 저장량 주장과 임의 탐색 기능이 양립하는지에 직접 영향을 준다.
RQ4가 이를 측정한다.}

\section{합성 팬텀}
\label{sec:phantom}

명제 \ref{prop:normal-error}, 보조정리 \ref{lem:quadratic}, 명제
\ref{prop:planar-disk}, \ref{prop:interp-eikonal}은 모두 \emph{수렴률 또는 정확한 항등식}에
대한 주장이다. 이를 확인하려면 정확한 참조가 필요하다. 즉 정확한 부호 거리 함수,
정확한 표면 법선, 그리고 $\spacing$을 임의로 바꿀 수 있는 능력이다. ACDC와 M\&Ms-2는
ED/ES label만 제공하며 어느 것도 제공하지 않는다.

따라서 해석적 참값을 갖는 합성 4D 팬텀을 구현하였다. 좌심실 내막을 수축하는 장축
회전 타원체로 모델링하고, 강체 변환이 거리를 보존하므로 기울임 아래에서도 부호 거리
함수가 \emph{정확}하게 유지된다. 비등방 표본화($h_z\gg h_x=h_y$), Rician 잡음, 매끄러운
명암 불균일, 그리고 혈액 풀 안의 심근 밝기 papillary muscle 덩어리를 포함한다.
papillary muscle은 Chan--Vese의 조각마다 상수 가정을 \emph{의도적으로} 위반하며, 정답
표면에는 포함되지 않는다(ACDC 관례와 일치). 또한 혈액과 심근 밝기가 주기에 따라
맥동하도록 하였다. 그렇지 않으면 외형 잔차가 항등적으로 0이 되어 RQ3가 공허해진다.

\caveat{%
팬텀은 \textbf{검증 도구이며 데이터셋이 아니다}. 실제 해부학, scanner 변이, 질환에
대한 강건성에 대해 아무것도 말해 주지 않는다. 실제 데이터 경로는 별도로 구현되어
있으나 실행되지 않았다.}

\section{검증 상태}
\label{sec:impl-status}

\subsection{실행되지 않았다}

본 구현은 PyTorch를 설치할 수 없는 환경에서 작성되었다(패키지 색인 접근 불가, 캐시된
wheel 없음, GPU 없음). 따라서 \textbf{PyTorch를 필요로 하는 어떤 부분도 실행된 적이
없다}. 확인된 것은 다음이다.

\begin{itemize}[leftmargin=2em]
  \item 49개 모듈과 4개 테스트 파일의 문법 (\texttt{compileall} 통과),
  \item 모든 상대 import가 정의된 이름으로 해소됨,
  \item \texttt{\_\_all\_\_} 항목과 모든 테스트 import의 해소,
  \item 텐서 형상과 수식의 행별 검토. 이 과정에서 실제 결함 4건을 발견하고 수정하였다.
  \footnote{(i) 0/1 텐서에 대한 \texttt{argmax}가 첫 최댓값을 반환한다는 보장이 없어
  광선 행진의 첫 교차 탐색이 비결정적이었던 문제, (ii) 밀도 제어로 surfel 수가 변할 때
  잔차 배열의 형상이 어긋나던 문제, (iii) 작은 격자에서 팬텀 타원체가 볼륨 경계에서
  잘려 닫힌 표면 가정이 깨지던 문제, (iv) 정사영 카메라에 원근 보정 보간을 잘못
  적용하던 문제.}
\end{itemize}

\subsection{실행된 검증: 이산 논리 커널}
\label{sec:kernel-verification}

PyTorch를 필요로 하지 않는 부분이 있다. 이산 색인 연산과 참조표이다. 이 영역은
오프바이원 하나가 예외를 발생시키지 않고 결과를 조용히 오염시키므로 위험도가 높다.
이를 표준 라이브러리만으로 재구현하여 독립 참조와 대조하였다. 41개 항목이
\emph{실제로 실행되어} 통과하였다.

\begin{table}[htbp]
\centering
\small
\caption{실행된 커널 검증. 본 논문에서 실제로 실행된 유일한 검증이다.}
\label{tab:kernel-checks}
\begin{tabular}{@{}p{6.2cm}p{7.2cm}@{}}
\toprule
검사 항목 & 독립 참조 \\
\midrule
큐브의 6면체 분해가 정확히 타일링 & 부피 합, 점 포함 횟수 \\
marching tetrahedra 16개 case의 위상 정확성 & 교차 간선 집합과의 일치 \\
(surfel, 타일) 페어 열거 & 브루트포스 이중 루프 \\
타일별 깊이 정렬과 절단 & 최근접 $M$개 보존 \\
점유 비트마스크 팩/언팩 & 8의 모든 나머지 길이에서 왕복 \\
타원체 부호 거리 이분법 브래킷 & 해석적 구, 브루트포스 최근접점 \\
삼선형 가중치의 합과 순서 & 아핀 정확성, 코너 대응 \\
공액경사법 수렴 & $n$회 내 정확해 \\
Godunov 상류 재초기화 스텐실 & $\norm{\nabla\lev}=1$이 고정점 \\
저장량 손익분기 대수 & 식 \eqref{eq:breakeven} 역산 시 $\CR=1$ \\
수축 프로파일의 연속성 & ES 및 주기 경계에서의 도약 \\
\bottomrule
\end{tabular}
\end{table}

\begin{remark}[검증이 검증 자체의 결함을 찾은 사례]
첫 실행에서 ``모든 내부 점이 정확히 하나의 사면체에 속한다''는 항목이 216개 표본 중
96개에서 실패하였다. 6개 사면체가 공유하는 내부 평면은 $x=y$, $y=z$, $x=z$ 세 개이며,
격자점 $(i,j,k)$, $i,j,k\in\{1,\dots,6\}$ 중 이 평면 위에 놓인 개수를 포함--배제로
계산하면
\[
  3\cdot36 - 3\cdot6 + 6 = 96
\]
으로 \emph{정확히 일치}하였다. 즉 분해는 옳고 검사가 사면체 경계면 위의 점을
표본화한 것이 원인이었다. 표본을 무리수 오프셋으로 옮기고, 경계면 위의 점이 정확히
두 사면체에 속하는지를 별도 항목으로 추가하여 해결하였다. 오류 개수가 예측값과
정확히 맞았다는 사실 자체가 분해의 정확성에 대한 강한 증거였다.
\end{remark}

\subsection{단계별 스모크 진단}

실행되지 않은 구현의 첫 실행에서는 여러 결함이 동시에 드러날 것이 예상된다. 첫 실패에서
멈추는 테스트 도구는 ``한 번 실행에 결함 하나'' 루프를 강요한다. 따라서 21개 독립
스테이지로 구성된 스모크 진단을 함께 구현하였다. 각 스테이지가 자신의 예외를 잡으므로
한 번 실행이 모든 실패를 보고하고, 의존성 순서로 배치되어 첫 실패가 이후 실패를 설명한다.
선행 스테이지가 실패한 스테이지는 \texttt{SKIP}으로 표시하여 무관한 실패와 구분한다.

스테이지는 실행되지 않은 구현에서 가장 발생 가능성이 높은 결함을 겨냥한다. 부호 규약
역전, 프레임 직교성 파괴, 수반 연산자 오류, 재생이 정준 집합을 변형하는 문제 등이며,
오류 메시지가 추정 원인을 명시한다. 예를 들어 구 곡률이 $-2/R$ 대신 $+2/R$로 측정되면
``내부/외부 규약이 역전되었음을 뜻한다''를 함께 출력한다.

\section{재현성}

전역 난수 상태에 의존하지 않도록 팬텀 표본 추출과 정준 초기화에 생성기를 명시적으로
전달한다. 면적 균일 표본 추출은 CPU에서 수행한 뒤 색인만 장치로 옮기므로 계산 장치와
무관하게 비트 단위로 재현된다. 모든 결과 파일은 장치, PyTorch 버전, GPU 이름을 포함한
환경 기록을 함께 저장한다. 타이밍 주장은 하드웨어 의존적이므로 이 기록 없이는 의미가
없다.
